\section{Round-sixteen positive closure: finite-time observability, process Gaussianity, and quadratic recovery}
\label{sec:r16-b3}

The covariance convention is fixed once: the isonormal collision measure has
control \(A_f\), and the biased intensity appears through the single factor
\(\sqrt q\).  Finite-time observability, process tightness, and second
Mosco convergence are proved as separate theorems rather than inferred from a
first-order LDP.

\subsection{Weighted balance complex}

Let \(m>8\).  Put
\[
 X=C([0,T];H^{-s}_{x,v}(1+|v|^m)),
 \qquad
 Y=L^2(A_f),
\]
with \(s\) larger than the phase-space dimension divided by two plus two.  A
smooth cotangent pair is \((p,\psi)\), where \(p\) has transport derivative and
endpoint traces in the weighted dual and \(\psi\in Y\).  Define
\[
 \mathfrak D(p,\psi)=\Delta p+\psi,
 \qquad
 \mathfrak G r=(r,-\Delta r).
\]
The primal linearized balance operator \(\mathfrak B\) maps density/contact
increments to the endpoint and transport residual.

\begin{theorem}[Finite-time kinetic observability and closed range]
\label{thm:r16-b3-observability}
On a sufficiently small regular biased phase chart around a Maxwellian
short-time solution, there is \(C_T<\infty\) such that
\[
 \inf_{\chi\in\mathcal N}
 \|p-\chi\|_{\mathcal P_T}^2
 \le C_T\left(
 \int_0^T\!\int q|\Delta p|^2dA_f
 +\|(-\partial_t-v\cdot\nabla_x-\mathcal L_{f,q}^*)p
   \|_{L^2_tH^{-1}_{x,v}}^2
 +\|p_T\|_{L^2(f_T)}^2
 \right),
\]
where \(\mathcal N\) is the finite collision-invariant space and
\(\mathcal P_T\) includes endpoint, transport, velocity-gradient, and trace
norms.  Consequently \(\mathfrak B\) has closed range and the complete
annihilator of balanced increments is
\(\operatorname{Ran}\mathfrak G\) together with the normalized collision
invariants.
\end{theorem}

\begin{proof}
Decompose \(p\) into the five hydrodynamic collision invariants and their
orthogonal complement in the Maxwellian velocity space.  The symmetrized
linearized collision operator has a spectral gap on the microscopic
component.  Commuting free transport with one spatial derivative transfers
microscopic coercivity to the nonconstant hydrodynamic modes.  Integrating the
backward equation against the standard hypocoercive cross term
\(\langle\nabla_x p,\nabla_v p\rangle\) and choosing its coefficient small
gives
\[
 {d\over dt}\mathcal E(p)+c\|p\|_{H^1_{x,v}/\mathcal N}^2
 \le C\|(-\partial_t-v\cdot\nabla_x-\mathcal L^*)p\|_{H^{-1}}^2.
\]
Endpoint integration controls the constant spatial modes.  The Green trace
theorem of B2 identifies the boundary contribution with the collision
Dirichlet form and controls its weighted trace.  Smooth regular biased
coefficients are a bounded perturbation of the Maxwellian operator; a
Gronwall argument preserves coercivity on the declared source/time chart.
The estimate is exactly the adjoint observability inequality for
\(\mathfrak B\).  The Hilbert closed-range theorem gives closed range and the
annihilator identity.  Normalization fixes the finite invariant directions.
\end{proof}

\begin{theorem}[Integral duality and exact gauge]
\label{thm:r16-b3-duality}
For every finite-moment pair \((f,\Gamma)\),
\[
 \sup_{p,\psi}\left\{
 \langle\mathfrak B(f,\Gamma),p\rangle
 +\int\psi\,d\Gamma
 -\int(e^{\Delta p+\psi}-1)dA_f
 \right\}
\]
is infinite unless the pair is balanced; on the balanced set it equals
\[
 \int\ell(d\Gamma/dA_f)dA_f.
\]
The identifiable cotangent is the quotient class represented by
\(z=\Delta p+\psi\), and for a positive intensity \(q\) the optimal class is
\(z=\log q\).
\end{theorem}

\begin{proof}
If balance fails, closed range from
Theorem~\ref{thm:r16-b3-observability} separates the residual by a cotangent;
choosing \(\psi=-\Delta p\) and scaling makes the supremum infinite.  On the
balanced set the first term becomes \(\int\Delta p\,d\Gamma\), so the
functional depends only on \(z=\Delta p+\psi\).  The scalar conjugacy of
\(e^z-1\) and \(q\log q-q+1\), followed by monotone truncation in the
exponential Orlicz heart, gives the entropy.  Exactness of
\(r\mapsto(r,-\Delta r)\mapsto\Delta p+\psi\) gives the quotient.
\end{proof}

\subsection{Covariance constructed from finite cumulants}

At a regular exposed source, let \((f,qA_f)\) be the biased law.  Let
\(\mathbb W\) be the isonormal Gaussian random measure with control measure
\(A_f\).  Define the collision martingale by
\[
 M_t(p,\psi)=
 \int_{[0,t]}\sqrt q\,(\Delta p+\psi)\,d\mathbb W.
\]
Then
\[
 \mathbb E M_t(p,\psi)M_t(\tilde p,\tilde\psi)
 =\int_{[0,t]}q(\Delta p+\psi)
   (\Delta\tilde p+\tilde\psi)dA_f.
\]

\begin{theorem}[Joint process-valued covariance]
\label{thm:r16-b3-covariance}
There is a unique centered Gaussian process
\((\zeta,\Xi)\in C([0,T];H^{-s})\times H^{-s}_{\rm coll}\) satisfying the
linearized biased balance
\[
 d\langle\zeta_t,\phi\rangle
 =\langle\zeta_t,
 v\cdot\nabla_x\phi+\mathcal L_{f,q}^*\phi\rangle dt
 +dM_t(\phi,0),
\]
and the joint contact relation
\[
 \Xi(\psi)=M_T(0,\psi)
 +\int\psi\,D(qA_f)[\zeta].
\]
Its initial covariance is the B1 microcanonical Schur complement.  For every
pair of density/contact tests its covariance is
\[
 \Sigma((p,\psi),(\tilde p,\tilde\psi))
 =\Sigma_0(p_0,\tilde p_0)
 +\int q(\Delta p+\psi)(\Delta\tilde p+\tilde\psi)dA_f,
\]
where \(p,\tilde p\) solve the corresponding backward equations.  This
continuous bilinear form defines a Radon Gaussian measure on the nuclear
path-space triple.
\end{theorem}

\begin{proof}
The energy estimate of Theorem~\ref{thm:r16-b3-observability} gives a bounded
backward solution operator from terminal/contact tests to \(\mathcal P_T\).
The displayed covariance is positive and continuous because it is the sum of
the positive initial form and the \(L^2(A_f)\) collision isometry.  Minlos'
theorem on the chosen nuclear test space constructs the Gaussian Radon law.
Variation of constants gives the density process, and differentiating the
biased contact intensity gives the second relation.  Uniqueness follows by the
same backward duality.  Pure contact tests have variance
\(\int q\psi^2dA_f\), and density--contact cross terms are the mixed factors in
the single product above.
\end{proof}

\subsection{Uniform interval cumulants and tightness}

For a smooth test pair supported in a deterministic interval \(I\), let
\(Z_\varepsilon(I)\) be the centered
\(\sqrt{\mu_\varepsilon}\)-scaled density/contact field.  The B2 connected
trajectory expansion localizes in time.

\begin{lemma}[Localized cumulant density]
\label{lem:r16-b3-cumulants}
For every \(k\ge2\), uniformly over intervals \(I\subset[0,T]\),
\[
 |\operatorname{cum}_k Z_\varepsilon(I)|
 \le C_k\mu_\varepsilon^{1-k/2}
 \left(|I|+{1\over\mu_\varepsilon}\right).
\]
For \(k=2\), the limit is the covariance measure of
Theorem~\ref{thm:r16-b3-covariance}; the term
\(\mu_\varepsilon^{-1}\) is the microscopic diagonal contribution.
Furthermore
\[
 \mathbb P\left(
 \sup_{I:\,|I|\le h}|Z_\varepsilon(I)|>\eta
 \right)
 \le C\eta^{-4}(h^2+h/\mu_\varepsilon)+Ce^{-c\eta\sqrt\mu}.
\]
\end{lemma}

\begin{proof}
A connected cluster contributing to an interval-supported cumulant must have
one distinguished contact or trajectory root in \(I\).  Integrating its root
time gives \(|I|\); when two observations hit the same microscopic contact,
the diagonal cell contributes \(\mu^{-1}\).  The fixed-horizon factorial
majorant of B2 controls all remaining labels uniformly.  Scaling gives the
first estimate.  The fourth-moment expansion into connected cumulants gives
\(C(h^2+h/\mu)\) on each grid cell.  A dyadic chaining union bound, together
with the exponential contact-count moment for the maximal single jump, gives
the supremum estimate.
\end{proof}

\begin{theorem}[Prepared joint Gaussian process limit]
\label{thm:r16-b3-clt}
Under the B1-conditioned hard-sphere law,
\[
 \sqrt{\mu_\varepsilon}
 \bigl(\pi^\varepsilon-f,\Gamma^\varepsilon-qA_f\bigr)
 \Longrightarrow(\zeta,\Xi)
\]
in the weighted distribution path topology.  The convergence is joint with
all finite source likelihood ratios on the regular chart.
\end{theorem}

\begin{proof}
Normal source convergence makes finite-dimensional cumulants of order at
least three vanish and identifies the second cumulant with
Theorem~\ref{thm:r16-b3-covariance}.  Lemma~\ref{lem:r16-b3-cumulants} gives
Aldous--Mitoma tightness on a countable nuclear determining family and controls
the maximal jump.  Every limit solves the unique Gaussian martingale problem.
Analyticity on a larger real source ball gives a uniform
\(L^{1+\delta}\) bound for centered likelihood ratios, hence joint
convergence.
\end{proof}

\subsection{Second epi-differentiability}

Let \(I\) be the B2 rate at the exposed path \(x_*=(f,qA_f)\).  For a tangent
\(z\), define
\[
 d_e^2I(x_*|\theta)(z)
 =\Gamma\hbox{-}\lim_{h\downarrow0}
 {2[I(x_*+hz)-I(x_*)-h\langle\theta,z\rangle]\over h^2}.
\]

\begin{theorem}[Quadratic recovery and Mosco limit]
\label{thm:r16-b3-mosco}
The second epi-derivative exists on the balanced tangent space and equals
\[
 d_e^2I(x_*|\theta)(z)=
 \begin{cases}
 \|\Sigma^{-1/2}z\|^2,
 &z\in\operatorname{Ran}\Sigma^{1/2},\\
 +\infty,&\text{otherwise}.
 \end{cases}
\]
For every \(z\) in the Cameron--Martin range there are smooth positive
balanced pairs \(x_h\) with
\(h^{-1}(x_h-x_*)\to z\) and
\[
 I(x_h)=I(x_*)+h\langle\theta,z\rangle
 +{h^2\over2}\|\Sigma^{-1/2}z\|^2+o(h^2).
\]
\end{theorem}

\begin{proof}
Write \(z=\Sigma^{1/2}y\) and approximate \(y\) by finite source directions
\(y_m\).  Exponential tilting by \(hy_m\) produces a smooth positive biased
path \(x_{h,m}\).  Third-derivative Cauchy bounds for the B1/B2 pressure give,
uniformly for fixed \(m\), the exact quadratic expansion of both mean and
Legendre cost with remainder \(O(h^3\|y_m\|^3)\).  Choose
\(m=m(h)\uparrow\infty\) slowly so that the source approximation and cubic
remainder are \(o(h^2)\).  This is the required quadratic recovery sequence.

For the liminf, Fenchel's inequality with sources
\(\theta+hy_m\), followed by \(m\to\infty\), gives the closed quadratic form
on \(\operatorname{Ran}\Sigma^{1/2}\).  If \(z\) is outside that range, the
supremum over finite source directions diverges.  Strong recovery and weak
lower semicontinuity give Mosco convergence.
\end{proof}

\begin{theorem}[Round-sixteen B3 closure]
\label{thm:r16-b3-main}
The balance operator has a proved finite-time closed range, the collision
Gaussian uses one biased-intensity factor, the complete initial/dynamic
covariance is a Radon process law, deterministic-interval cumulants give
process tightness, and second epi-differentiability has an explicit quadratic
recovery sequence.
\end{theorem}

\begin{proof}
Combine Theorems~\ref{thm:r16-b3-observability},
\ref{thm:r16-b3-duality}, \ref{thm:r16-b3-covariance},
\ref{thm:r16-b3-clt}, and \ref{thm:r16-b3-mosco}, and
Lemma~\ref{lem:r16-b3-cumulants}.
\end{proof}
